(1) 传播相位分别为 \(\phi_x=2\pi n_xd/\lambda\)、\(\phi_y=2\pi n_yd/\lambda\)。去掉共同相位 \(e^{i\phi_x}\)，输出 Jones 矢量
\[
\boldsymbol E_{\rm out}=E_0(\cos\alpha,e^{i\delta}\sin\alpha),\qquad
\delta=\frac{2\pi(n_y-n_x)d}{\lambda}.
\]

(2) 当 \(\alpha=\pi/4\)，两轴振幅相等。若 \(\delta=\pm\pi/2\pmod{2\pi}\)，实电场两分量相位差为四分之一周期，其端点在 xy 平面绕成圆，且总强度恒定；对应的两种符号为相反旋向的圆偏振。

(3) 轴角 \(\beta \) 的理想线检偏器投影矢量为 \((\cos\beta,\sin\beta)\)，通过后的电场幅
\[
E_\beta=E_0[\cos\alpha\cos\beta+e^{i\delta}\sin\alpha\sin\beta].
\]
故
\[
I_\beta=I_0[\cos^2\alpha\cos^2\beta+
\sin^2\alpha\sin^2\beta+
2\cos\delta\sin\alpha\cos\alpha\sin\beta\cos\beta].
\]
这里 \(I_0\) 为入射总强度。

(4) 若 \(\delta=\pi\)，输出与 \(E_0(\cos\alpha,-\sin\alpha)\) 等价，是与 x 轴夹角 \(-\alpha\) 的线偏振。因输入角为 \(+\alpha\)，偏振方向镜像反转，波片主轴若旋转也会使输出偏振方向按相应几何关系转动。
